\documentclass[11pt,a4paper]{report}
\usepackage[utf8]{inputenc}
\usepackage[T1]{fontenc}
\usepackage{lmodern}
\usepackage[spanish,es-tabla]{babel}
\usepackage{geometry}
\geometry{margin=2.5cm}
\usepackage{hyperref}
\usepackage{xcolor}
\usepackage{listings}
\usepackage{graphicx}
\usepackage{booktabs}
\usepackage{longtable}
\usepackage{enumitem}
\usepackage{titlesec}
\usepackage{fancyhdr}
\usepackage{tabularx}
\usepackage{multirow}

% --- Hyperlink colors ---
\hypersetup{
  colorlinks=true,
  linkcolor=blue!70!black,
  urlcolor=blue!70!black,
  citecolor=blue!70!black,
}

% --- Code listings ---
\definecolor{codebg}{rgb}{0.95,0.95,0.98}
\lstset{
  backgroundcolor=\color{codebg},
  basicstyle=\ttfamily\small,
  breaklines=true,
  frame=single,
  framesep=4pt,
  keywordstyle=\color{blue!70!black},
  commentstyle=\color{gray},
  stringstyle=\color{red!60!black},
  numbers=left,
  numberstyle=\tiny\color{gray},
  tabsize=2,
  language=TypeScript,
}

% --- Section formatting ---
\titleformat{\chapter}[display]
  {\normalfont\huge\bfseries\color{blue!70!black}}{\chaptertitlename\ \thechapter}{20pt}{\Huge}
\titleformat{\section}
  {\normalfont\Large\bfseries\color{blue!60!black}}{\thesection}{1em}{}
\titleformat{\subsection}
  {\normalfont\large\bfseries}{\thesubsection}{1em}{}

% --- Header/footer ---
\pagestyle{fancy}
\fancyhf{}
\fancyhead[L]{\leftmark}
\fancyhead[R]{Arbiter Selector — Documentación}
\fancyfoot[C]{\thepage}

% --- Title ---
\title{\vspace{2cm}\Huge\textbf{Arbiter Selector}\\[0.5cm]
  \Large Documentación del Proyecto\\[0.3cm]
  \normalsize TFT Set 17 — Spec-Driven Development}
\author{Equipo de Desarrollo}
\date{\today}

\begin{document}

\maketitle
\thispagestyle{empty}
\cleardoublepage

% --- Table of Contents ---
\tableofcontents
\cleardoublepage

% ============================================================
\chapter{Resumen Ejecutivo}
% ============================================================

\section{Propósito del Proyecto}

\textbf{Arbiter Selector} es una aplicación web interactiva de una sola página (SPA) que permite a los jugadores de \textit{Teamfight Tactics} (TFT) Set 17 explorar, simular y compartir combinaciones del rasgo \textbf{Arbiter} (Sentencia).

El rasgo Arbiter consiste en una matriz de $9 \times 9$ combinaciones de \textit{causa-efecto} (``leyes divinas''), donde no todos los emparejamientos son válidos. Esta herramienta cubre tres modos de interacción:

\begin{enumerate}
  \item \textbf{Simulador}: Selección aleatoria de $3 \times 3$ leyes imitando la experiencia del juego, con función de re-roll y combinaciones válidas garantizadas.
  \item \textbf{Explorador}: Selección libre de causa/efecto con acceso a la matriz completa, filtrado de efectos válido por causa, y URLs compartibles.
  \item \textbf{Base de Datos}: Tabla de referencia completa $9 \times 9$ y tarjetas de referencia de campeones.
\end{enumerate}

\section{Stack Tecnológico}

\begin{table}[h]
\centering
\begin{tabular}{@{}ll@{}}
\toprule
\textbf{Capa} & \textbf{Tecnología} \\
\midrule
Build & Vite 6 \\
Framework & React 19 + TypeScript 5.7 \\
Estilos & CSS vanilla con custom properties \\
Iconos UI & lucide-react \\
Iconos stats & SVG inline (9 componentes) \\
Testing & Vitest (26 tests) \\
Deploy & GitHub Pages (GitHub Actions) \\
\bottomrule
\end{tabular}
\end{table}

\section{Métricas del Proyecto}

\begin{itemize}
  \item \textbf{Tasks completadas}: 31 de 31 (100\%)
  \item \textbf{Fases}: 6 (Foundation, Core Components, Utilities, Tab Implementation, App Assembly, Testing)
  \item \textbf{Archivos creados}: $\sim$54
  \item \textbf{Tests}: 26/26 passing
  \item \textbf{Bundle}: 212 KB JS, 12 KB CSS (68 KB gzipped total)
  \item \textbf{TypeScript}: 0 errores en modo strict
\end{itemize}

\cleardoublepage

% ============================================================
\chapter{Exploración}
% ============================================================

\section{Estado Actual}

Proyecto greenfield. No existe código previo. Los requisitos se definieron a partir de los datos oficiales de la wiki de TFT Set 17.

\section{Enfoques Considerados}

\subsection{Estructura de Datos}

\begin{table}[h]
\centering
\begin{tabularx}{\textwidth}{@{}lXll@{}}
\toprule
\textbf{Enfoque} & \textbf{Descripción} & \textbf{Pros} & \textbf{Contras} \\
\midrule
A: Matriz plana & \texttt{Record<CauseId, Record<EffectId, LawValues | null>>} & O(1) lookup, type-safe, simple & Difícil de extender para otros traits \\
B: Arrays normalizados & Arrays separados + Set de pares válidos & Extensible, separación clara & Más archivos, lookups complejos \\
C: Sistema genérico & Interfaz \texttt{TraitDefinition} abstracta & Future-proof & Over-engineering para un solo trait \\
\bottomrule
\end{tabularx}
\end{table}

\textbf{Decisión}: Enfoque A (matriz plana). La app es específicamente para Arbiter. Si se agregan más traits después, se puede refactorizar.

\subsection{Codificación de URL}

\begin{table}[h]
\centering
\begin{tabular}{@{}lXll@{}}
\toprule
\textbf{Opción} & \textbf{Ejemplo} & \textbf{Pros} & \textbf{Contras} \\
\midrule
A: Slugs legibles & \texttt{?cause=attacks-3x\&effect=max-health} & Debuggable, compartible, SEO-friendly & URLs más largas \\
B: IDs numéricos & \texttt{?c=1\&e=3} & Más cortas & No legibles, frágiles \\
C: Base64 & \texttt{?law=eyJ...} & Compacto & Opaco, difícil de debuggear \\
\bottomrule
\end{tabular}
\end{table}

\textbf{Decisión}: Opción A (slugs legibles). Solo 2 parámetros necesarios. Los slugs funcionan como documentación.

\subsection{Gestión de Estado}

\textbf{Decisión}: React Context. El estado es pequeño (1 tab enum, 3 causas + 3 efectos para simulador, 1 causa + 1 efecto para explorador), las actualizaciones son infrecuentes, y no hay estado derivado complejo. No se necesita Zustand/Redux.

\section{Riesgos Identificados}

\begin{enumerate}
  \item \textbf{Datos incompletos}: Algunos valores de la wiki son rangos (ej. ``8\%/9\%''). Se muestran tal cual con comentarios de fuente.
  \item \textbf{Selección aleatoria}: Si la matriz es muy dispersa, los re-rolls pueden ser frecuentes. Solución: pre-computar pares válidos y muestrear de esos.
  \item \textbf{Responsive mobile}: La tabla $9 \times 9$ es muy ancha. Scroll horizontal con headers sticky o layout basado en cards.
  \item \textbf{Iconos}: Se necesitan íconos para stats. Se usan SVGs inline para cero dependencias externas.
  \item \textbf{Parches de balance}: Los valores pueden cambiar entre parches. Capa de datos versionada por parche.
\end{enumerate}

\cleardoublepage

% ============================================================
\chapter{Propuesta}
% ============================================================

\section{Intención}

El rasgo Arbiter en TFT Set 17 es el más complejo del set — una matriz de $9 \times 9$ de combinaciones causa-efecto donde no todas son válidas. Los jugadores necesitan una herramienta interactiva para entender, explorar y compartir combinaciones válidas de leyes. Esta app cubre ese vacío.

\section{Alcance}

\subsection{Dentro del Alcance}
\begin{itemize}
  \item Tab Simulador: Selección aleatoria $3 \times 3$ con re-roll y combinaciones válidas garantizadas
  \item Tab Explorador: Selección libre con acceso a la matriz completa, filtrado de efectos, URLs compartibles
  \item Tab Base de Datos: Tabla de referencia $9 \times 9$, tarjetas de campeones
  \item URLs compartibles con query params legibles
  \item Diseño responsive (desktop, tablet, mobile)
  \item Tema oscuro con estética Space Gods (púrpura/cósmico)
\end{itemize}

\subsection{Fuera del Alcance}
\begin{itemize}
  \item Simulación de combate o cálculos de daño
  \item Team builder o planificación de composiciones
  \item Otros traits de TFT (solo Arbiter)
  \item Backend o API — app estática 100\% client-side
  \item Cuentas de usuario o persistencia más allá de URLs
\end{itemize}

\section{Criterios de Éxito}

\begin{enumerate}
  \item La app se deploya a GitHub Pages y carga sin errores
  \item Los 3 tabs son completamente funcionales
  \item Las URLs son compartibles y restauran el estado correctamente
  \item Los datos coinciden con los valores de la wiki
  \item La tabla muestra las 81 celdas con indicadores válidos/inválidos correctos
  \item El simulador garantiza al menos una combinación válida por roll
  \item El layout mobile es usable (probado a 375px)
  \item TypeScript strict mode pasa con cero errores
  \item Todos los tests pasan (vitest)
\end{enumerate}

\cleardoublepage

% ============================================================
\chapter{Especificación}
% ============================================================

\section{Dominio 1: Capa de Datos}

\subsection{Tipos de Causa (REQ-DL-001 a REQ-DL-003)}

El sistema define exactamente 9 identificadores de causa:

\begin{table}[h]
\centering
\small
\begin{tabular}{@{}ll@{}}
\toprule
\textbf{ID} & \textbf{Descripción} \\
\midrule
\texttt{attacks-3x} & Unidad ataca 3 veces \\
\texttt{combat-start-per-star} & Inicio de combate: por nivel de estrella \\
\texttt{enemy-dies} & Muere una unidad enemiga \\
\texttt{deals-damage-10x} & Unidad causa daño 10 veces \\
\texttt{combat-start-per-interest} & Inicio de combate: por interés ganado \\
\texttt{every-4s} & Cada 4 segundos en combate \\
\texttt{spends-50-mana} & Unidad gasta 50 de maná \\
\texttt{combat-start-if-rerolled} & Inicio de combate: si se hizo reroll \\
\texttt{first-time-below-40hp} & Primera vez bajo 40\% de HP \\
\bottomrule
\end{tabular}
\end{table}

\subsection{Tipos de Efecto (REQ-DL-004 a REQ-DL-006)}

El sistema define exactamente 9 identificadores de efecto:

\begin{table}[h]
\centering
\small
\begin{tabular}{@{}ll@{}}
\toprule
\textbf{ID} & \textbf{Descripción} \\
\midrule
\texttt{max-health} & Vida máxima \\
\texttt{shield} & Escudo \\
\texttt{armor-mr} & Armadura y Resistencia Mágica \\
\texttt{attack-speed} & Velocidad de Ataque \\
\texttt{ability-power} & Poder de Habilidad \\
\texttt{mana} & Maná \\
\texttt{chance-gold} & Probabilidad de oro \\
\texttt{chance-reroll} & Probabilidad de reroll \\
\texttt{chance-leona} & Probabilidad de Leona \\
\bottomrule
\end{tabular}
\end{table}

\subsection{Matriz de Leyes (REQ-DL-007 a REQ-DL-011)}

\begin{itemize}
  \item Estructura: mapeo $9 \times 9$ donde cada par causa-efecto mapea a \texttt{LawValues} o \texttt{null}
  \item Formato de valores: \texttt{tier2} (string) y \texttt{tier3} (string)
  \item Combinaciones inválidas: representadas como \texttt{null}
  \item Completitud: 81 entradas totales (72 válidas, 9 inválidas)
  \item Performance: lookups en O(1)
\end{itemize}

\subsection{Datos de Campeones (REQ-DL-012 a REQ-DL-014)}

\begin{table}[h]
\centering
\begin{tabular}{@{}llll@{}}
\toprule
\textbf{Nombre} & \textbf{Costo} & \textbf{Sinergia} & \textbf{Rol} \\
\midrule
Leona & 1g & Vanguard & Tank \\
Zoe & 2g & Conduit & Support \\
Diana & 3g & Challenger & Carry \\
LeBlanc & 4g & Shepherd & Carry \\
\bottomrule
\end{tabular}
\end{table}

\section{Dominio 2: Compartición por URL}

\begin{itemize}
  \item \textbf{Formato}: \texttt{?cause=attacks-3x\&effect=max-health} (slugs legibles, kebab-case)
  \item \textbf{Tab}: \texttt{?tab=simulator|explorer|database}
  \item \textbf{Restauración}: Al cargar la app, se parsea la URL y se restaura el estado
  \item \textbf{Validación}: Slugs inválidos o combinaciones null se ignoran gracefulmente
  \item \textbf{Copy-to-clipboard}: Botón con feedback visual (check verde por 2 segundos)
  \item \textbf{Prohibido}: No se usa Base64 ni IDs numéricos en URLs
\end{itemize}

\section{Dominio 3: Simulador}

\begin{itemize}
  \item Selecciona exactamente 3 causas y 3 efectos aleatorios
  \item \textbf{Garantía}: al menos una combinación válida existe entre las $3 \times 3$ opciones
  \item Algoritmo: pick 3 causas $\rightarrow$ unión de efectos válidos $\rightarrow$ pick 3 efectos
  \item Sin duplicados en causas ni efectos
  \item Re-roll independiente (no influenciado por rolls anteriores)
  \item Máximo 10 intentos antes de fallback determinista
\end{itemize}

\section{Dominio 4: Componentes UI}

\subsection{Navegación por Tabs (REQ-UI-001 a REQ-UI-005)}
\begin{itemize}
  \item 3 tabs: Simulator, Explorer, Database
  \item Tab activo indicado visualmente
  \item URL se sincroniza al cambiar de tab
  \item Tab por defecto: Simulator
\end{itemize}

\subsection{Selectores de Causa y Efecto (REQ-UI-006 a REQ-UI-017)}
\begin{itemize}
  \item Cards seleccionables con hover y selected states
  \item Efectos filtrados por causa seleccionada (solo válidos)
  \item Toggle: clickear la misma card la deselecciona
\end{itemize}

\subsection{Display de Resultados (REQ-UI-018 a REQ-UI-022)}
\begin{itemize}
  \item Muestra tier-2 y tier-3 lado a lado
  \item Diferenciación visual entre tiers
  \item Mensaje apropiado para combinaciones inválidas
\end{itemize}

\subsection{Tabla de Base de Datos (REQ-UI-023 a REQ-UI-028)}
\begin{itemize}
  \item Matriz completa $9 \times 9$ con headers sticky
  \item Celdas válidas muestran valores tier2 $\rightarrow$ tier3
  \item Celdas inválidas muestran ``---'' con opacidad reducida
  \item Zebra striping para legibilidad
  \item Scroll horizontal en mobile
\end{itemize}

\subsection{Tema Oscuro (REQ-UI-032 a REQ-UI-036)}
\begin{itemize}
  \item Paleta basada en hue 250-270 (púrpura/cósmico)
  \item Todos los colores como CSS custom properties
  \item Contraste WCAG AA mínimo
\end{itemize}

\section{Requisitos No Funcionales}

\begin{table}[h]
\centering
\small
\begin{tabular}{@{}ll@{}}
\toprule
\textbf{ID} & \textbf{Requisito} \\
\midrule
NFR-001 & TypeScript strict mode, cero errores \\
NFR-002 & Sin librerías externas de estado (solo React Context) \\
NFR-003 & CSS vanilla con custom properties, sin CSS-in-JS \\
NFR-004 & Cero dependencias de assets externos (SVGs inline) \\
NFR-005 & 100\% client-side, sin backend \\
NFR-006 & Vite como build tool \\
NFR-007 & Vitest para unit tests \\
NFR-008 & Carga en menos de 2 segundos \\
NFR-009 & Bundle bajo 100 KB gzipped \\
\bottomrule
\end{tabular}
\end{table}

\cleardoublepage

% ============================================================
\chapter{Diseño Técnico}
% ============================================================

\section{Enfoque Técnico}

La app se organiza alrededor de un componente \texttt{App} que posee el estado del tab y la sincronización con URL. Cada tab es un módulo autocontenido con sus sub-componentes. Los datos fluyen unidireccionalmente desde \texttt{arbiter-data.ts} a través de React Context hacia los componentes hoja.

\subsection{Principios Core}
\begin{itemize}
  \item \textbf{Única fuente de verdad}: \texttt{arbiter-data.ts} es el único archivo de datos
  \item \textbf{URL como espejo de estado}: Cada acción del usuario actualiza la URL; el reload restaura el estado
  \item \textbf{Context slices separados}: \texttt{TabContext}, \texttt{SimulatorContext}, \texttt{ExplorerContext} — separados para evitar re-renders innecesarios
  \item \textbf{Capa de utilidades puras}: Toda la lógica de leyes vive en \texttt{lib/} como funciones puras, completamente testeables
\end{itemize}

\section{Decisiones de Arquitectura}

\subsection{Estructura de Datos: Matriz Plana}

\textbf{Decisión}: \texttt{Record<CauseId, Record<EffectId, LawValues | null>>}

\textbf{Racional}:
\begin{itemize}
  \item Lookup O(1): \texttt{matrix[causeId][effectId]} — sin iteración
  \item Auto-documentado: las keys son slugs semánticos
  \item TypeScript proporciona autocompletado en las 9 causas y 9 efectos
  \item Pares inválidos son explícitamente \texttt{null}
\end{itemize}

\subsection{Gestión de Estado: React Context}

Tres contextos separados para evitar re-renders innecesarios:
\begin{itemize}
  \item \texttt{TabProvider} — \texttt{activeTab}, \texttt{setActiveTab}
  \item \texttt{SimulatorProvider} — \texttt{causes[]}, \texttt{effects[]}, \texttt{roll()}
  \item \texttt{ExplorerProvider} — \texttt{selectedCause}, \texttt{selectedEffect}
\end{itemize}

\subsection{Codificación de URL: Slugs Legibles}

\texttt{URLSearchParams} con slugs kebab-case. Sin Base64, sin IDs numéricos.

\subsection{Organización de Componentes: Feature Folders}

\begin{lstlisting}[language=,basicstyle=\ttfamily\footnotesize]
src/components/
├── layout/          — Header, TabNav
├── simulator/       — SimulatorTab, LawSelector, LawGrid
├── explorer/        — ExplorerTab, CausePicker, EffectPicker
├── database/        — DatabaseTab, MatrixTable, ChampionCards
└── shared/          — LawResult, CopyButton, EmptyState
\end{lstlisting}

\subsection{Estilos: CSS Custom Properties + Vanilla CSS}

Paleta Space Gods (hue 260):
\begin{itemize}
  \item \texttt{--surface-0} a \texttt{--surface-3} — capas de fondo
  \item \texttt{--text-primary}, \texttt{--text-secondary}, \texttt{--text-tertiary} — jerarquía de texto
  \item \texttt{--accent} (hsl 275, 70\%, 62\%) — color principal
  \item Escala de spacing de 8px (\texttt{--space-1} a \texttt{--space-9})
\end{itemize}

\subsection{Estrategia de Iconos}

Dos niveles:
\begin{itemize}
  \item \textbf{Iconos de stats} (HP, Shield, AS, AP, Mana, Gold, etc.): SVGs inline en \texttt{src/icons/}
  \item \textbf{Iconos de UI} (flechas, refresh, copy): \texttt{lucide-react} con tree-shaking
\end{itemize}

\section{Flujo de Datos}

\begin{lstlisting}[language=,basicstyle=\ttfamily\footnotesize]
arbiter-data.ts (matriz 9x9, campeones)
    │
    ▼
lib/ utilities (law-utils.ts, law-encoding.ts)
    │
    ▼
React Context Layer (TabContext, SimulatorContext, ExplorerContext)
    │
    ▼
Component Tree (App → Header → TabNav → Tabs)
    │
    ▼
URL Synchronization (decodeUrl on mount, encodeUrl on change)
\end{lstlisting}

\section{Estructura Final de Archivos}

\begin{lstlisting}[language=,basicstyle=\ttfamily\footnotesize]
arbiter-selector/
├── .github/workflows/deploy.yml
├── package.json
├── vite.config.ts
├── vitest.config.ts
├── tsconfig.json / tsconfig.app.json / tsconfig.node.json
├── index.html
├── .gitignore
├── public/favicon.svg
└── src/
    ├── main.tsx
    ├── App.tsx
    ├── index.css (~500 líneas)
    ├── types.ts
    ├── context/
    │   ├── index.ts
    │   ├── AppProvider.tsx
    │   ├── SimulatorProvider.tsx
    │   └── ExplorerProvider.tsx
    ├── data/
    │   ├── arbiter-data.ts (matriz 9x9, 72 válidas, 9 null)
    │   └── constants.ts
    ├── lib/
    │   ├── index.ts
    │   ├── law-utils.ts (7 funciones)
    │   ├── law-utils.test.ts (12 tests)
    │   ├── law-encoding.ts (8 funciones)
    │   └── law-encoding.test.ts (14 tests)
    ├── icons/ (9 componentes SVG)
    └── components/
        ├── layout/ (Header, TabNav)
        ├── shared/ (LawResult, CopyButton, EmptyState)
        ├── simulator/ (LawSelector, LawGrid, SimulatorTab)
        ├── explorer/ (CausePicker, EffectPicker, ExplorerTab)
        └── database/ (MatrixTable, ChampionCards, DatabaseTab)
\end{lstlisting}

\cleardoublepage

% ============================================================
\chapter{Implementación — Tasks}
% ============================================================

\section{Fase 1: Foundation / Infrastructure (6 tasks)}

\begin{table}[h]
\centering
\small
\begin{tabular}{@{}lll@{}}
\toprule
\textbf{Task} & \textbf{Descripción} & \textbf{Archivos} \\
\midrule
1.1 & Project scaffolding (Vite + React + TS) & package.json, tsconfigs, vite.config.ts, index.html \\
1.2 & TypeScript types & src/types.ts \\
1.3 & Data layer (matriz 9x9) & src/data/arbiter-data.ts \\
1.4 & Constants (labels, arrays) & src/data/constants.ts \\
1.5 & CSS foundation (Space Gods theme) & src/index.css \\
1.6 & Stat icons (9 SVGs) & src/icons/*.tsx \\
\bottomrule
\end{tabular}
\end{table}

\section{Fase 2: Core Components (7 tasks)}

\begin{table}[h]
\centering
\small
\begin{tabular}{@{}lll@{}}
\toprule
\textbf{Task} & \textbf{Descripción} & \textbf{Archivos} \\
\midrule
2.1 & AppProvider con TabContext + URL sync & src/context/AppProvider.tsx \\
2.2 & Header component & src/components/layout/Header.tsx \\
2.3 & TabNav component & src/components/layout/TabNav.tsx \\
2.4 & LawResult component & src/components/shared/LawResult.tsx \\
2.5 & CopyButton component & src/components/shared/CopyButton.tsx \\
2.6 & EmptyState component & src/components/shared/EmptyState.tsx \\
2.7 & Barrel exports & src/context/index.ts, etc. \\
\bottomrule
\end{tabular}
\end{table}

\section{Fase 3: Utilities (3 tasks)}

\begin{table}[h]
\centering
\small
\begin{tabular}{@{}lll@{}}
\toprule
\textbf{Task} & \textbf{Descripción} & \textbf{Archivos} \\
\midrule
3.1 & law-utils (7 funciones puras) & src/lib/law-utils.ts \\
3.2 & law-encoding (8 funciones URL) & src/lib/law-encoding.ts \\
3.3 & Barrel export & src/lib/index.ts \\
\bottomrule
\end{tabular}
\end{table}

\section{Fase 4: Tab Implementation (11 tasks)}

\begin{table}[h]
\centering
\small
\begin{tabular}{@{}lll@{}}
\toprule
\textbf{Task} & \textbf{Descripción} & \textbf{Archivos} \\
\midrule
4.1 & SimulatorProvider & src/context/SimulatorProvider.tsx \\
4.2 & LawSelector & src/components/simulator/LawSelector.tsx \\
4.3 & LawGrid (layout horizontal) & src/components/simulator/LawGrid.tsx \\
4.4 & SimulatorTab & src/components/simulator/SimulatorTab.tsx \\
4.5 & ExplorerProvider + URL sync & src/context/ExplorerProvider.tsx \\
4.6 & CausePicker (9 causas) & src/components/explorer/CausePicker.tsx \\
4.7 & EffectPicker (filtrado) & src/components/explorer/EffectPicker.tsx \\
4.8 & ExplorerTab & src/components/explorer/ExplorerTab.tsx \\
4.9 & MatrixTable (9x9) & src/components/database/MatrixTable.tsx \\
4.10 & ChampionCards (4 campeones) & src/components/database/ChampionCards.tsx \\
4.11 & DatabaseTab & src/components/database/DatabaseTab.tsx \\
\bottomrule
\end{tabular}
\end{table}

\section{Fase 5: App Assembly (3 tasks)}

\begin{table}[h]
\centering
\small
\begin{tabular}{@{}lll@{}}
\toprule
\textbf{Task} & \textbf{Descripción} & \textbf{Archivos} \\
\midrule
5.1 & App.tsx con tab routing & src/App.tsx \\
5.2 & main.tsx con TabProvider & src/main.tsx \\
5.3 & GitHub Pages deploy workflow & .github/workflows/deploy.yml \\
\bottomrule
\end{tabular}
\end{table}

\section{Fase 6: Testing (3 tasks)}

\begin{table}[h]
\centering
\small
\begin{tabular}{@{}lll@{}}
\toprule
\textbf{Task} & \textbf{Descripción} & \textbf{Resultado} \\
\midrule
6.1 & Unit tests law-utils & 12 tests passing \\
6.2 & Unit tests law-encoding & 14 tests passing \\
6.3 & Build verification & tsc: 0 errores, build: success \\
\bottomrule
\end{tabular}
\end{table}

\section{Descubrimiento Clave}

Vite y Vitest shippean versiones incompatibles del bundler interno. \textbf{No se puede mergear} la configuración de Vitest en \texttt{vite.config.ts} — requiere un \texttt{vitest.config.ts} separado con entorno jsdom.

\cleardoublepage

% ============================================================
\chapter{Resultados de Verificación}
% ============================================================

\section{Build Final}

\begin{table}[h]
\centering
\begin{tabular}{@{}ll@{}}
\toprule
\textbf{Métrica} & \textbf{Valor} \\
\midrule
TypeScript (tsc --noEmit) & 0 errores \\
Build (npm run build) & Success (7.97s) \\
HTML & 0.63 KB \\
CSS & 12.10 KB (2.49 KB gzipped) \\
JS & 212.18 KB (65.56 KB gzipped) \\
\textbf{Total gzipped} & \textbf{$\sim$68 KB} (bajo límite de 100 KB) \\
Tests (vitest run) & 26/26 passing \\
\bottomrule
\end{tabular}
\end{table}

\section{Estado del Proyecto}

\textbf{TODAS LAS 31 TASKS COMPLETADAS} — El proyecto es completamente funcional, testeado y listo para deployment.

\end{document}
